\documentclass[10pt,a4paper]{article}
\usepackage[utf8]{inputenc}
\usepackage[ngerman]{babel}
\usepackage{amsmath}
\usepackage{amsfonts}
\usepackage{amssymb}
\usepackage{mathtools}
\usepackage{microtype}
\usepackage{authblk}
\usepackage{titlesec}

% section title formating
\renewcommand{\thesection}{\Roman{section}}
\titleformat{\section}[block]{\normalsize\bfseries}{\thesection.}{1em}{\MakeUppercase}
\titlespacing*{\section}{0pt}{11pt plus 2pt minus 2pt}{11pt plus 2pt minus 2pt}

\renewcommand{\thesubsection}{\arabic{section}.\Alph{subsection}}
\titleformat{\subsection}{\normalsize\bfseries}{\thesubsection}{1em}{}
\titlespacing*{\subsection}{0pt}{11pt plus 2pt minus 2pt}{11pt plus 2pt minus 2pt}

\renewcommand{\thesubsubsection}{\thesubsection.\arabic{subsubsection}}
\titleformat{\subsubsection}{\normalsize\itshape}{\thesubsubsection}{1em}{}
\titlespacing*{\subsubsection}{0pt}{11pt plus 2pt minus 2pt}{11pt plus 2pt minus 2pt}
   
% custom macros
\newcommand{\mydef}[1]{\textls{#1}}
\newcommand{\nul}{\textrm{---}}
\newcommand{\trans}{^\mathsf{T}}

% %%%%%%%
% title %
% %%%%%%%

\title{Herleitung der Eulerschen Kreiselgleichungen aus den Lagrange-Gleichungen auf~$\mathrm{SO}(3)$}
\author{Hendrik Meier}

\affil{Immenstaad am Bodensee}
\renewcommand\Affilfont{\itshape}

\date{November 2023}


% %%%%%%%%%%
% document %
% %%%%%%%%%%

\begin{document}
\maketitle

\section{Motivation}

\subsection{Eulersche Gleichungen}

Wirken keine äußeren Momente, so ist der Drehimpuls~$L^I$, aufgefasst als Vektor im $\mathbb{R}^3$, im Inertialsystem ($I$) erhalten, $\dot{L}^I=0$.
Im körperfesten Bezugssystem ($B$) ist der Drehimpuls~$L^B=R\trans L^I$ durch die Zeitabhängigkeit der Lagematrix~$R\in\mathrm{SO}(3)$ im Allgemeinen nicht konstant, sondern genügt der durch einfaches Differenzieren erhaltenen Gleichung
\begin{align}
	\dot{L}^B &= L^B \times \omega
	\ ,
\end{align}
wobei die (körperfest gemessene) Winkelgeschwindigkeit~$\omega$ kinematisch durch die  Lage gegeben ist, $[\omega\times]=\Omega=R\trans\dot{R}$.

Körperfester Drehimpuls~$L^B$ und Winkelgeschwindigkeit~$\omega$ sind über den Trägheitstensor~$J$, hier aufgefasst als Operator auf $\mathbb{R}^3$, miteinander verknüpft,
\begin{align}
\label{eq:l_b}
	L^B &= J\cdot\omega
	\ .
\end{align}
Wir erhalten somit
\begin{align}
\label{eq:euler_homogen}
	J\cdot\dot{\omega} &=  -\omega\times(J\cdot\omega)
	\ .
\end{align}
Dies sind die \mydef{Eulerschen Kreiselgleichungen} für den Fall, dass keine äußeren Momente wirken.

\subsection{Kinetische Energie}

Ein \mydef{starrer Körper} sei definiert durch $n\geq 1$ Massepunkte der Massen~$m_i$, $i\in\{1,\ldots,n\}$, die bezüglich eines gewählten körperfesten Ursprungs~$O$ die zeitlich konstanten körperfesten Positionen~$r_i^B\in\mathbb{R}^3$ einnehmen.
Wir nehmen an, dass der Ursprung~$O$ des körperfesten Systems mit jenem des betrachteten Inertialsystems zu jeder Zeit übereinstimmt, betrachten also ausschließlich rotatorische Bewegung des starren Körpers.

Die kinetische Energie des starren Körpers lautet im Inertialsystem
\begin{align}
	T = \frac{1}{2}\sum_i m_i 
	\left|
		\dot{r}_i^I
	\right|^2
	\ .
\end{align}
Durch Übergang ins körperfeste System erhalten wir
\begin{align}
	T &= 
	\frac{1}{2}\sum_i m_i 
	\left|
		\frac{d}{dt}
		\left(
			Rr_i^B
		\right)
	\right|^2
	=
	\frac{1}{2}\sum_i m_i 
	\left|
		\omega\times r_i^B
	\right|^2
	\nonumber\\
	&=
	\frac{1}{2}\sum_i m_i
	\left(
	\left|r_i^B\right|^2
	\left|\omega\right|^2	
	- 
	\langle
		r_i^B,\omega
	\rangle^2
	\right)
	\nonumber\\
	&=
	\frac{1}{2} J\cdot(\omega\otimes\omega)
	\ ,
\end{align}
wobei wir den \mydef{Trägheitstensor}~$J$ hier in natürlicher Weise als $(0,2)$-Tensor
\begin{align}
\label{eq:j_02}
	J &=
	\sum_i m_i
	\left(
	\left|r_i^B\right|^2
	\langle\nul,\nul\rangle
	- 
	\langle
		r_i^B,\nul
	\rangle
	\otimes
	\langle
		r_i^B,\nul
	\rangle
	\right)	
\end{align}
definieren.
Drücken wir die kinetische Energie durch die (kartesischen) Koordinaten der Winkelgeschwindigkeit $\omega=(\omega_\alpha)_{\alpha\in\{1,2,3\}}$ aus,
so erhalten wir
\begin{align}
\label{eq:t_omega}
	T &= \frac{1}{2} \sum_{\alpha,\beta=1}^3 J_{\alpha\beta}\omega_\alpha\omega_\beta
	\ .
\end{align}
Die kinetische Energie nach Gl.~(\ref{eq:t_omega}) hängt nur von der (körperfest gemessenen) Winkelgeschwindigkeit~$\omega$ ab, wohingegen die Lage~$R$ in der Formel nicht erscheint.
Dies ist physikalisch auf Grund der Isotopie des Raumes verständlich.

Fasst man hingegen Gl.~(\ref{eq:t_omega}) ``in naiver Weise'' als Lagrange-Funktion auf, so scheint es unmöglich, die Eulerschen Gleichungen~(\ref{eq:euler_homogen}) zu erhalten.
Diesen scheinbaren Widerspruch aufzulösen, ist Ziel dieses Aufsatzes.

\section{Lagrange-Gleichungen des starren Körpers}

Wir betrachten einen Keim der Bewegung des starren Körpers.
Diesen setzen wir über einem den Zeitpunkt~$t=0$ umfassenden Intervall~$I$ an,
\begin{align}
\label{eq:r_t}
	R: I &\rightarrow \mathrm{SO}(3)
	\ ,\\\nonumber
	t &\mapsto R(t)
	\ .
\end{align}
Für die kinetische Energie gilt
\begin{align}
	T &= 
	\frac{1}{2}\sum_i m_i 
	\left|				
		\dot{R}(t)r_i^B
	\right|^2
	=
	\frac{1}{2}\sum_i m_i 
	\left\langle
		r_i^B, \dot{R}\trans(t)\dot{R}(t)\,r_i^B
	\right\rangle
	\nonumber\\
\label{eq:t_rdot}
	&= \frac{1}{2}
	\mathrm{Sp}\left[
		\tilde{J}
		\cdot
		\dot{R}\trans(t)\dot{R}(t)
	\right]
\end{align}
mit dem Operator~$\tilde{J}$, der für $u\in\mathbb{R}^3$ durch
\begin{align}
	\tilde{J}\cdot u = \sum_i m_i 
	r_i^B
	\left\langle
		r_i^B, u
	\right\rangle		
\end{align}
definiert ist.
Bezeichnen wir mit~$J$ auch den Gl.~(\ref{eq:j_02}) entsprechenden $(1,1)$-Tensor, gelten offenbar die Beziehungen
\begin{align}
	J &= \mathrm{Sp}[\tilde{J}]\,\mathrm{id}_{\mathbb{R}^3} - \tilde{J}
	\ \ \text{bzw.}\ \ 
	\tilde{J} = 
	\frac{1}{2}
	\mathrm{Sp}[J]\,\mathrm{id}_{\mathbb{R}^3}
	- J
	\ .
\end{align}

\subsection{Generalisierte Koordinaten}

Ziel ist die Herleitung der Lagrange-Gleichungen zum Zeitpunkt~$t=0$.
Dazu führen wir \mydef{Normalkoordinaten}~$\theta$ für den Tangentialraum an~$R_0=R(0)$ ein.
Diesen identifizieren wir mit $\mathbb{R}^3$, wobei die Lie-Klammer durch das Standardvektorprodukt ausgedrückt wird.
Wir verwenden ferner~$\Theta$ als Bezeichnung für die $\theta$ entsprechende schiefsymmetrische Matrix, d.h. $\Theta\cdot u=\theta\times u$ für alle $u\in\mathbb{R}^3$.

Die Exponentialabbildung in die Lie-Gruppe~$\mathrm{SO}(3)$ nimmt in den betrachteten Normalkoordinaten die Form
\begin{align}
	\mathrm{Exp}: 
	\mathrm{T}_{R_0}\mathrm{SO}(3)\cong\mathbb{R}^3
	&\rightarrow \mathrm{SO}(3)
	\nonumber\\
	\theta
	&\mapsto
	R_0\,\exp\Theta
\end{align}
an, wobei $\exp$ die Matrixexponentialfunktion ist.
Für die schiefsymmetrischen Argumente~$\Theta$ können wir $\exp$ auch mit Hilfe der Formel von Rodrigues
\begin{align}
\label{eq:rodrigues}
	\exp\Theta
	&= \mathrm{id}_{\mathbb{R}^3} 
	+ \frac{\sin|\theta|}{|\theta|}\,\Theta
	+ \frac{1 - \cos|\theta|}{|\theta|^2}\,\Theta^2
\end{align}
auswerten.

Die Normalkoordinaten~$\theta$ dienen uns im Folgenden als die generalisierten Koordinaten zur Herleitung der Lagrange-Gleichungen.
Den Keim der Bewegung des starren Körpers aus Gl.~(\ref{eq:r_t}) betrachten wir entsprechend als Weg
\begin{align}
\label{eq:theta_t}
	t\mapsto\theta(t)
\end{align}
mit $\theta(0)=0$, so dass $\mathrm{Exp}[\theta(t)]=R(t)$ für alle $t\in I$.

In diesen Koordinaten entspricht die Geschwindigkeit~$\dot{\theta}(0)$ der körperfest gemessenen Winkelgeschwindigkeit~$\omega_0$ zum Zeitpunkt~$t=0$.
Es ist nämlich
\begin{align}
	\Omega_0 \coloneqq \Omega(0) &=
	\left.	
	R\trans(t)\frac{d}{dt}R(t)
	\right|_{t=0}
	=
	\left.	
	\frac{d}{dt}\exp\Theta(t)
	\right|_{t=0}
	\nonumber\\
	\label{eq:omega_dot_theta}
	&=
	\left.	
	\frac{d}{dt}(\mathrm{id}_{\mathbb{R}^3} + \Theta(t) + \ldots)
	\right|_{t=0}
	=\dot{\Theta}(0)
	\ ,
\end{align}
da $\Theta(0)=0$ ist.
Für die zum Zeitpunkt~$t=0$ körperfest gemessene Winkelbeschleunigung~$\dot{\omega}_0$ erhalten wir analog
\begin{align}
	\dot{\Omega}_0 \coloneqq \dot{\Omega}(0) &=
	\left.
	\frac{d}{dt}
	\left(
	R\trans(t)\frac{d}{dt}R(t)
	\right)
	\right|_{t=0}
	\nonumber\\
	&=
	\left.
	\frac{d}{dt}
	\left[
	(\mathrm{id}_{\mathbb{R}^3} - \Theta(t) + \ldots)
	\frac{d}{dt}
	\left(
		\mathrm{id}_{\mathbb{R}^3} 
		+ \Theta(t)
		+ \frac{1}{2}\Theta^2(t)
		+ \ldots
	\right)
	\right]
	\right|_{t=0}
	\nonumber\\
	\label{eq:dot_omega_ddot_theta}
	&=\ddot{\Theta}(0)
	\ .
\end{align}
Die körperfest gemessenen Winkelgeschwindigkeiten und -beschleunigungen stimmen damit zum Zeitpunkt $t=0$ mit den entsprechenden Ableitungen der Normalkoordinaten überein.

\subsection{Lagrange-Funktion}

In Abwesenheit von außen wirkender Momente ist die \mydef{Lagrange-Funktion}~$\mathcal{L}$ allein durch die kinetische Energie gegeben.
Einsetzen der generalisierten Koordinaten in Gl.~(\ref{eq:t_rdot}) liefert
\begin{align}
\label{eq:lagrange}
	\mathcal{L}(\theta, \dot{\theta})
	&=
	\frac{1}{2}
	\mathrm{Sp}\left[
		\tilde{J}
		\cdot
		\left(
			\frac{d}{dt}
			\exp[-\Theta]
		\right)		
		\left(
			\frac{d}{dt}
			\exp[\Theta]
		\right)	
	\right]
	\ .
\end{align}
Die Ableitungen sind darin rein formal aufzufassen, um die Abhängigkeit von $\theta$ und $\dot{\theta}$ darzustellen.
Eine explizite Darstellung von~$d\exp[\Theta]/dt$ durch die Formel~(\ref{eq:rodrigues}) von Rodrigues ergibt etwa
\begin{align}
	\frac{d}{dt}
	\exp[\Theta]
	&= 
	\frac{|\theta|\cos|\theta|-\sin|\theta|}{|\theta|^2}
	\frac{\langle
		\theta,\dot{\theta}
	\rangle}{|\theta|}
	\,\Theta
	+ \frac{\sin|\theta|}{|\theta|}\,\dot{\Theta}
	\nonumber\\
	&\quad
	+ \frac{|\theta|\sin|\theta|-2(1 - \cos|\theta|)}{|\theta|^3}
	\frac{\langle
		\theta,\dot{\theta}
	\rangle}{|\theta|}\,\Theta^2
	+ \frac{1 - \cos|\theta|}{|\theta|^2}\,
	\left(
		\Theta\dot{\Theta} + \dot{\Theta}\Theta
	\right)
	\ .
\end{align}
Die Lagrange-Funktion~$\mathcal{L}(\theta, \dot{\theta})$, Gl.~(\ref{eq:lagrange}), ist quadratisch in~$\dot{\theta}$.
Taylor-Entwicklung in $\theta$ liefert
\begin{align}
	\mathcal{L}(\theta, \dot{\theta})
	&=
	-\frac{1}{2}
	\mathrm{Sp}\left[
		\tilde{J}
		\cdot
		\left(
			\dot{\Theta}^2
			+ \frac{1}{2}
			\left(
				\Theta\dot{\Theta}^2
				-
				\dot{\Theta}^2\Theta
			\right)
			+ \mathrm{O}(|\dot{\theta}|^2|\theta|^2)
		\right)	
	\right]
	\ .
\end{align}
Die nicht ausgeschriebenen Terme~$\mathrm{O}(|\dot{\theta}|^2|\theta|^2)$ spielen für die abschließende Auswertung zum Zeitpunkt~$t=0$ keine Rolle, da mit $\theta=0$ auch deren Ableitung verschwindet.
Auf Grund der Symmetrie von $\tilde{J}$ folgt
\begin{align}
\label{eq:lagrange_approx}
	\mathcal{L}(\theta, \dot{\theta})
	&\simeq
	-\frac{1}{2}
	\mathrm{Sp}\left[
		\tilde{J}
		\cdot
		\left(
			\mathrm{id}_{\mathbb{R}^3}
			+
			\Theta
		\right)
		\cdot
		\dot{\Theta}^2
	\right]
\end{align}
unter Beibehaltung nur der später relevanten Terme.
Mit den bekannten Rechenregeln für Vektorprodukte lässt sich die Spur umschreiben als
\begin{align}
\label{eq:l_approx}
	\mathcal{L}(\theta, \dot{\theta})
	&\simeq
	-\frac{1}{2}\left(
		\langle
			\dot{\theta},
			\tilde{J}\cdot\dot{\theta}
		\rangle
		-
		|\dot{\theta}|^2\,
		\mathrm{Sp}\tilde{J}
		+
		\langle
		\theta\times\dot{\theta},
		\tilde{J}\cdot\dot{\theta}
		\rangle		
	\right)
	\nonumber\\
	&=
	\frac{1}{2}\left(
		\langle
			\dot{\theta},
			J\cdot\dot{\theta}
		\rangle
		-
		\langle
			\theta\times\dot{\theta},
			J\cdot\dot{\theta}
		\rangle
	\right)
	\ .
\end{align}
Gleichung~(\ref{eq:l_approx}) zeigt bereits, warum der ``naive'' Ansatz, Gl.~(\ref{eq:t_omega}) als Lagrange-Funktion zu verwenden, unvollständig ist:
Gleichung~(\ref{eq:t_omega}) für die kinetische Energie entspricht offenbar dem Wert von~$\mathcal{L}(\theta, \dot{\theta})$ an der Stelle~$\theta=0$, ohne aber darüber hinaus Informationen über die Abhängigkeit der kinetischen Energie von den Koordinaten~$\theta$ zu enthalten.


Für die Gradienten der Lagrange-Funktion nach~$\theta$ bzw. $\dot{\theta}$ verwenden wir die Bezeichnungen $\nabla_{\theta}\mathcal{L}$ bzw. $\nabla_{\dot{\theta}}\mathcal{L}$.
Zum Beispiel erhalten wir für den hier als Tangentialvektor aufgefassten zu~$\theta$ konjugierten Impuls
\begin{align}
\label{eq:l_conj}
	L &= \nabla_{\dot{\theta}}(\theta, \dot{\theta})
	\simeq
	J\cdot\dot{\theta}
	+
	\frac{1}{2}
	\left(
	\theta\times(J\cdot\dot{\theta})
	-J\cdot
	(\theta\times\dot{\theta})
	\right)
	\ .
\end{align}
Für den Zeitpunkt~$t=0$ des betrachteten, durch Gln.~(\ref{eq:r_t}) und~(\ref{eq:theta_t}) beschriebenen Weges in $\mathrm{SO}(3)$ folgt mit $\theta(0)=0$ sowie $\dot{\theta}(0)=\omega_0$, siehe Gl.~(\ref{eq:omega_dot_theta}), dass 
\begin{align}
	L(0) &= J\cdot\omega_0
\end{align}
gleich dem in Gl.(\ref{eq:l_b}) angegebenen Drehimpuls~$L^B$ im körperfesten System ist.

\subsection{Lagrange-Gleichungen}

In den (generalisierten) Koordinaten~$(\theta, \dot{\theta})$ der betrachteten lokalen Trivialisierung des Tangentialbündels~$\mathrm{T}\mathrm{SO}(3)$ lauten die Lagrange-Gleichungen
\begin{align}
	\frac{d}{dt}\nabla_{\dot{\theta}}\mathcal{L}
	=
	\nabla_{\theta}\mathcal{L}
	\ .
\end{align}
Einsetzen von Gl.~(\ref{eq:l_approx}) bzw. für die linke Seite direkt von Gl.~(\ref{eq:l_conj}) gibt
\begin{align}
	J\cdot\ddot{\theta}
	+
	\frac{1}{2}
	\left(
	\dot{\theta}\times(J\cdot\dot{\theta})
	+ \theta\times(J\cdot\ddot{\theta})
	-J\cdot
	(\theta\times\ddot{\theta})
	\right)
	&\simeq
	-\frac{1}{2}
	\dot{\theta}\times(J\cdot\dot{\theta})
	\ ,
\end{align}
also
\begin{align}
	J\cdot\ddot{\theta}
	+
	\dot{\theta}\times(J\cdot\dot{\theta})
	&\simeq
	\frac{1}{2}
	\left(
	J\cdot
	(\theta\times\ddot{\theta})
	- \theta\times(J\cdot\ddot{\theta})
	\right)
	\ .
\end{align}
Die Auswertung des betrachteten Weges in $\mathrm{SO}(3)$, vgl. Gln.~(\ref{eq:r_t}) und~(\ref{eq:theta_t}), zum Zeitpunkt~$t=0$ liefert mit $\theta(0)=0$ sowie $\dot{\theta}(0)=\omega_0$ und $\ddot{\theta}(0)=\dot{\omega}_0$, siehe Gln.~(\ref{eq:omega_dot_theta}) und~(\ref{eq:dot_omega_ddot_theta}), die Beziehung
\begin{align*}
	J\cdot\dot{\omega}_0
	+
	\omega_0\times(J\cdot\omega_0)
	&= 0
	\ .
\end{align*}
Dies sind wiederum die (homogenen) Eulerschen Kreiselgleichungen aus Gl.~(\ref{eq:euler_homogen}).


\end{document}